\documentclass[aspectratio=169,11pt]{beamer}
\usepackage{beamerucad}
\usepackage{listings}
\definecolor{ckw}{HTML}{2563AB}\definecolor{cst}{HTML}{2E8B57}\definecolor{ccm}{HTML}{8A8F98}
\definecolor{outbg}{HTML}{F4F5F7}\definecolor{outrule}{HTML}{2E8B57}
\definecolor{errbg}{HTML}{FBEDEC}\definecolor{errrule}{HTML}{C0392B}
\lstdefinestyle{py}{language=Python,basicstyle=\ttfamily\scriptsize,
 keywordstyle=\color{ckw}\bfseries,stringstyle=\color{cst},commentstyle=\color{ccm}\itshape,
 showstringspaces=false,columns=fullflexible,keepspaces=true,
 backgroundcolor=\color{ucadband!8},frame=leftline,framerule=2.5pt,rulecolor=\color{ucadband},
 xleftmargin=8pt,aboveskip=3pt,belowskip=1pt,basewidth=0.5em}
\lstdefinestyle{out}{basicstyle=\ttfamily\scriptsize\color{black!82},
 backgroundcolor=\color{outbg},frame=leftline,framerule=2.5pt,rulecolor=\color{outrule},
 xleftmargin=8pt,aboveskip=2pt,belowskip=1pt,columns=fullflexible,keepspaces=true}
\lstdefinestyle{err}{basicstyle=\ttfamily\scriptsize\color{black!82},
 backgroundcolor=\color{errbg},frame=leftline,framerule=2.5pt,rulecolor=\color{errrule},
 xleftmargin=8pt,aboveskip=2pt,belowskip=1pt,columns=fullflexible,keepspaces=true}
\lstset{style=py}
\newcommand{\resultat}{{\scriptsize\textbf{R\'esultat}~$\rightarrow$}\par\vspace{1pt}}

\title{Programmation Python — Séance 4}
\subtitle{Fonctions, modules et robustesse}
\author{Dr. El Hadji Bassirou TOURÉ}
\institute{Département de Mathématiques et Informatique\\ Faculté des Sciences et Techniques\\ Université Cheikh Anta Diop de Dakar}
\date{}
\setucadfoot{Programmation Python — Séance 4}
\begin{document}

\begin{frame}[plain,noframenumbering]\titlepage\end{frame}

% =====================================================================
\begin{frame}{Objectifs de la séance}
\begin{ucaddef}[Contenu]
{\small À mesure qu'un programme grandit, on répète du code et on rencontre des erreurs. Cette séance apprend à \textbf{organiser} (des fonctions réutilisables), à \textbf{réutiliser} (des modules tout prêts), à \textbf{conserver} (des fichiers) et à \textbf{fiabiliser} (gérer les erreurs).}
\end{ucaddef}
\begin{itemize}\small
  \item \textbf{Fonctions} : \texttt{def}, \texttt{return}, paramètres, valeurs par défaut, docstrings, portée.
  \item \textbf{Modules} : \texttt{import} pour réutiliser la bibliothèque standard.
  \item \textbf{Fichiers} : lire et écrire des données sur le disque.
  \item \textbf{Exceptions} : \texttt{try} / \texttt{except} pour ne pas planter.
\end{itemize}
{\small L'entraînement se fait dans le \textbf{notebook}, avec une cellule « À vous » après chaque notion.}
\end{frame}

\begin{frame}{Anatomie d'un appel de fonction}
\figslide{0.9}{figs/f4_fonction}{Les arguments fournis à l'appel remplissent les paramètres ; le corps s'exécute ; \texttt{return} renvoie une valeur.}
\end{frame}

% #####################################################################
\section{Les fonctions}

\begin{frame}{L'idée : un bloc de code réutilisable}
\begin{ucaddef}[L'idée en une phrase]
{\small Une \textbf{fonction} est un bloc nommé qu'on définit une fois avec \texttt{def}, puis qu'on \textbf{appelle} autant de fois que nécessaire. \texttt{return} renvoie un résultat que l'on peut réutiliser.}
\end{ucaddef}
{\small Le bénéfice est double : on \textbf{évite de répéter} le même code (une correction se fait à un seul endroit), et on \textbf{donne un nom} à une opération, ce qui rend le programme lisible. Le corps de la fonction est un bloc indenté après les deux points.}
\end{frame}

\begin{frame}[fragile]{Définir et appeler}
\begin{lstlisting}
def moyenne(a, b):
    return (a + b) / 2

print(moyenne(12, 16))
print(moyenne(9, 18))
\end{lstlisting}
\resultat
\begin{lstlisting}[style=out]
14.0
13.5
\end{lstlisting}
{\small \texttt{a} et \texttt{b} sont les \textbf{paramètres} ; \texttt{12} et \texttt{16} les \textbf{arguments} de l'appel. La même fonction sert pour n'importe quelles valeurs.}
\end{frame}

\begin{frame}[fragile]{Plusieurs paramètres}
\begin{lstlisting}
def message(prenom, ville):
    return f"{prenom} habite a {ville}."

print(message("Fatou", "Dakar"))
\end{lstlisting}
\resultat
\begin{lstlisting}[style=out]
Fatou habite a Dakar.
\end{lstlisting}
{\small Les arguments sont associés aux paramètres \textbf{dans l'ordre} : \texttt{prenom} reçoit « Fatou », \texttt{ville} reçoit « Dakar ».}
\end{frame}

\begin{frame}[fragile]{Valeurs par défaut}
\begin{ucadex}[Un paramètre peut avoir une valeur par défaut]
{\small Si l'argument n'est pas fourni, la valeur par défaut est utilisée.}
\end{ucadex}
\begin{lstlisting}
def saluer(nom, politesse="Bonjour"):
    return f"{politesse} {nom}"

print(saluer("Awa"))
print(saluer("Awa", "Bonsoir"))
\end{lstlisting}
\resultat
\begin{lstlisting}[style=out]
Bonjour Awa
Bonsoir Awa
\end{lstlisting}
\end{frame}

\begin{frame}[fragile]{Documenter : la docstring}
\begin{lstlisting}
def carre(x):
    "Renvoie le carre de x."
    return x * x

print(carre(5))
print(carre.__doc__)
\end{lstlisting}
\resultat
\begin{lstlisting}[style=out]
25
Renvoie le carre de x.
\end{lstlisting}
{\small La \textbf{docstring} est une courte phrase placée juste sous le \texttt{def}. Elle explique ce que fait la fonction, sans qu'on ait à lire son code.}
\end{frame}

\begin{frame}[fragile]{Piège — la portée d'une variable}
\begin{ucadpiege}[Une variable créée dans une fonction est locale]
{\small \texttt{resultat} n'existe que dans \texttt{calcul}. Y accéder à l'extérieur échoue.}
\end{ucadpiege}
\begin{lstlisting}
def calcul():
    resultat = 5 * 3

calcul()
print(resultat)
\end{lstlisting}
\begin{lstlisting}[style=err]
NameError: name 'resultat' is not defined
\end{lstlisting}
{\small Pour récupérer une valeur calculée dans une fonction, il faut la \textbf{renvoyer} avec \texttt{return}.}
\end{frame}

\begin{frame}[fragile]{\texttt{return} n'est pas \texttt{print}}
\begin{lstlisting}
def somme(a, b):
    return a + b

def affiche_somme(a, b):
    print(a + b)

x = somme(10, 5) + 100
print("x =", x)
y = affiche_somme(10, 5)
print("y =", y)
\end{lstlisting}
\resultat
\begin{lstlisting}[style=out]
x = 115
15
y = None
\end{lstlisting}
{\small \texttt{return} \textbf{renvoie} une valeur (réutilisable dans un calcul) ; une fonction qui se contente d'afficher renvoie \texttt{None}.}
\end{frame}

% #####################################################################
\section{Les modules}

\begin{frame}{L'idée : réutiliser du code tout prêt}
\begin{ucaddef}[L'idée en une phrase]
{\small Un \textbf{module} est une bibliothèque de code déjà écrit, qu'on \textbf{importe} pour l'utiliser. Python en fournit une riche collection en standard — inutile de tout réécrire.}
\end{ucaddef}
{\small Le module \texttt{math} donne les outils mathématiques, \texttt{random} le hasard, \texttt{datetime} les dates, etc. On importe une fois, en haut du fichier, puis on utilise. C'est un premier pas vers l'écosystème \texttt{numpy} / \texttt{pandas} qu'on découvrira à partir de la Séance 6.}
\end{frame}

\begin{frame}[fragile]{Importer un module : \texttt{import}}
\begin{lstlisting}
import math

print(math.sqrt(16))
print(math.pi)
\end{lstlisting}
\resultat
\begin{lstlisting}[style=out]
4.0
3.141592653589793
\end{lstlisting}
{\small Après \texttt{import math}, on accède à ses outils en les préfixant par \texttt{math.} — ce qui indique clairement d'où ils viennent.}
\end{frame}

\begin{frame}[fragile]{Importer un seul outil : \texttt{from ... import}}
\begin{lstlisting}
from math import sqrt

print(sqrt(25))
\end{lstlisting}
\resultat
\begin{lstlisting}[style=out]
5.0
\end{lstlisting}
{\small On importe \textbf{seulement} ce dont on a besoin ; \texttt{sqrt} s'utilise alors directement, sans préfixe. Pratique pour quelques fonctions bien identifiées.}
\end{frame}

% #####################################################################
\section{Les fichiers}

\begin{frame}{L'idée : conserver des données sur le disque}
\begin{ucaddef}[L'idée en une phrase]
{\small Les variables disparaissent à la fin du programme. Un \textbf{fichier} permet de \textbf{sauvegarder} des données et de les retrouver plus tard. On l'ouvre avec \texttt{open()} dans un bloc \texttt{with}, qui le \textbf{referme automatiquement}.}
\end{ucaddef}
{\small Le \textbf{mode} précise l'intention : \texttt{"w"} écrit (et écrase le contenu existant), \texttt{"a"} ajoute à la fin, \texttt{"r"} lit (c'est le mode par défaut). Un fichier texte est une simple suite de lignes, séparées par \texttt{\textbackslash n}.}
\end{frame}

\begin{frame}[fragile]{Écrire dans un fichier}
\begin{lstlisting}
with open("villes.txt", "w") as f:
    f.write("Dakar\n")
    f.write("Louga\n")
    f.write("Mbour\n")

print("fichier ecrit")
\end{lstlisting}
\resultat
\begin{lstlisting}[style=out]
fichier ecrit
\end{lstlisting}
{\small Le bloc \texttt{with} ouvre \texttt{villes.txt} en écriture, y écrit trois lignes, puis le referme tout seul en sortant du bloc. Le \texttt{\textbackslash n} marque la fin de chaque ligne.}
\end{frame}

\begin{frame}[fragile]{Lire un fichier}
{\small Tout le contenu d'un coup, avec \texttt{.read()} :}
\begin{lstlisting}
with open("villes.txt") as f:
    contenu = f.read()
print(contenu)
\end{lstlisting}
{\small Ou \textbf{ligne par ligne}, avec une boucle \texttt{for} (le plus courant) :}
\begin{lstlisting}
with open("villes.txt") as f:
    for ligne in f:
        print(ligne.strip())
\end{lstlisting}
\resultat
\begin{lstlisting}[style=out]
Dakar
Louga
Mbour
\end{lstlisting}
\end{frame}

\begin{frame}{Piège — le mode d'ouverture}
\begin{ucadpiege}[\texttt{"w"} écrase tout]
{\small Ouvrir un fichier existant en mode \texttt{"w"} \textbf{efface son contenu} avant d'écrire. Pour \textbf{compléter} un fichier sans le vider, on utilise le mode \texttt{"a"} (append). Pour seulement le lire, \texttt{"r"} (ou rien, c'est le défaut) — et écrire dans un fichier ouvert en lecture lève une erreur.}
\end{ucadpiege}
{\small Toujours privilégier le bloc \texttt{with} : il garantit que le fichier est refermé, même si une erreur survient au milieu.}
\end{frame}

% #####################################################################
\section{Les exceptions}

\begin{frame}[fragile]{L'idée : rattraper une erreur au lieu de planter}
\begin{ucaddef}[L'idée en une phrase]
{\small Certaines opérations \textbf{échouent} à l'exécution. Sans précaution, le programme s'arrête net :}
\end{ucaddef}
\begin{lstlisting}
nombre = int("abc")
\end{lstlisting}
\begin{lstlisting}[style=err]
ValueError: invalid literal for int() with base 10: 'abc'
\end{lstlisting}
{\small On \textbf{essaie} le code risqué dans un bloc \texttt{try}, et on \textbf{rattrape} l'erreur dans un bloc \texttt{except} — le programme continue au lieu de s'interrompre.}
\end{frame}

\begin{frame}[fragile]{\texttt{try} / \texttt{except}}
\begin{lstlisting}
try:
    nombre = int("abc")
    print(nombre)
except ValueError:
    print("Ce n'est pas un nombre valide.")
\end{lstlisting}
\resultat
\begin{lstlisting}[style=out]
Ce n'est pas un nombre valide.
\end{lstlisting}
{\small Le code du \texttt{try} s'exécute ; dès qu'une \texttt{ValueError} survient, on saute au \texttt{except}, qui affiche un message clair. Le programme n'a pas planté.}
\end{frame}

\begin{frame}[fragile]{Une fonction robuste}
\begin{ucadex}[Renvoyer une valeur de repli plutôt que planter]
{\small On précise \textbf{quelle} erreur on rattrape (\texttt{ZeroDivisionError}).}
\end{ucadex}
\begin{lstlisting}
def division_sure(a, b):
    try:
        return a / b
    except ZeroDivisionError:
        return "impossible"

print(division_sure(10, 2))
print(division_sure(10, 0))
\end{lstlisting}
\resultat
\begin{lstlisting}[style=out]
5.0
impossible
\end{lstlisting}
\end{frame}

\begin{frame}{Piège — ne pas tout masquer}
\begin{ucadpiege}[Rattraper l'erreur attendue, pas toutes]
{\small On indique l'erreur précise qu'on sait gérer (\texttt{ValueError}, \texttt{ZeroDivisionError}, \texttt{FileNotFoundError}\ldots). Un \texttt{except} \textbf{sans type}, qui rattrape \emph{tout}, cache aussi les erreurs qu'on n'avait pas prévues (une faute de frappe dans un nom, par exemple) et rend les bugs très difficiles à trouver.}
\end{ucadpiege}
{\small Régle : rattraper le \textbf{moins} possible — seulement ce que l'on sait traiter — et laisser remonter le reste.}
\end{frame}

% #####################################################################
\section{Synthèse}

\begin{frame}{À retenir — Séance 4}
\begin{ucadret}[L'essentiel]
{\small
\textbf{Fonction} : \texttt{def nom(parametres):} + \texttt{return} ; appelée avec des \textbf{arguments} ; peut avoir des \textbf{valeurs par défaut} et une \textbf{docstring}.\\[2pt]
\textbf{Portée} : une variable définie dans une fonction est \textbf{locale} ; pour en sortir une valeur, on la \texttt{return}. \texttt{return} $\neq$ \texttt{print}.\\[2pt]
\textbf{Modules} : \texttt{import math} puis \texttt{math.sqrt}, ou \texttt{from math import sqrt}.\\[2pt]
\textbf{Fichiers} : \texttt{with open("f.txt", "w") as f:} pour écrire, \texttt{"r"} pour lire ; \texttt{.read()} ou boucle \texttt{for} ligne par ligne ; \texttt{"w"} écrase, \texttt{"a"} ajoute.\\[2pt]
\textbf{Exceptions} : \texttt{try:} le code risqué, \texttt{except} l'erreur \textbf{précise} ; ne pas tout masquer.}
\end{ucadret}
\end{frame}

\begin{frame}{Pièges fréquents}
\begin{itemize}\small
  \item Utiliser dehors une variable \textbf{locale} à une fonction : \texttt{NameError} — la \texttt{return}.
  \item Attendre une valeur d'une fonction qui fait \texttt{print} au lieu de \texttt{return} (elle renvoie \texttt{None}).
  \item Oublier les deux points \texttt{:} ou l'indentation du corps de la fonction.
  \item Ouvrir en mode \texttt{"w"} un fichier qu'on voulait compléter : son contenu est \textbf{effacé}.
  \item Rattraper \textbf{toutes} les erreurs avec un \texttt{except} nu : les vrais bugs deviennent invisibles.
\end{itemize}
\end{frame}

\begin{frame}{Prochaine séance}
\begin{ucaddef}[Séance 5 — Programmation orientée objet]
{\small On sait écrire des fonctions. La Séance 5 introduit les \textbf{objets} : regrouper des données \emph{et} les fonctions qui les manipulent dans une même \textbf{classe} (\texttt{class}, \texttt{\_\_init\_\_}, \texttt{self}, méthodes), puis \textbf{spécialiser} une classe par \textbf{héritage} et \textbf{polymorphisme}.}
\end{ucaddef}
{\small À faire d'ici là : refaire les cellules « À vous » du notebook de la séance.}
\end{frame}

\begin{frame}{Références}
\begin{itemize}\small
  \item Severance, C. — \emph{Python for Everybody}, 2016.
  \item Downey, A. — \emph{Think Python}, 2e éd., O'Reilly, 2015.
  \item Matthes, E. — \emph{Python Crash Course}, 3e éd., No Starch Press, 2023.
  \item Python Software Foundation — \emph{The Python Tutorial} (Functions, Modules, Errors and Exceptions), 2024.
\end{itemize}
\end{frame}

\end{document}
